% answers_short.tex —— 2007 年·春季上海卷 · 简短答案速查
% 由 tools/make_short_answers.py 自 answers.tex 的 \gksolhead/\gkanswer 提取，
% 并与 2007/2007_数学真题_春季上海_答案要点.md 逐题比对后生成，勿手改（改源 MD 后重跑）。

\begin{gknotice}
\textbf{使用说明}：本节只给最终答案便于快速核对，完整推导见解析版。答案为编写者独立推导，\textbf{非官方标准答案}。
\end{gknotice}

\gkpart{填空题}
\begin{itemize}[leftmargin=2.6em]
  \item[\textbf{1.}] $\dfrac23$。
  \item[\textbf{2.}] $q=2$。
  \item[\textbf{3.}] $a=4$。
  \item[\textbf{4.}] $\pi$。
  \item[\textbf{5.}] $-3$。
  \item[\textbf{6.}] $5$。
  \item[\textbf{7.}] $2$。
  \item[\textbf{8.}] $\dfrac{3\pi}{4}$，即 $135^\circ$。
  \item[\textbf{9.}] $-1$。
  \item[\textbf{10.}] $120$ 人。
  \item[\textbf{11.}] 反函数为 $f^{-1}(x)=\begin{cases}\sqrt{x-1},&x\ge1,\\ \dfrac2x,&x<0.\end{cases}$，定义域为 $(-\infty,0)\cup[1,+\infty)$。
\end{itemize}

\gkpart{选择题}
\begin{center}\begin{tabular}{ccccc}
\toprule
题号 & 12 & 13 & 14 & 15 \\
\midrule
答案 & A & B & B & D \\
\bottomrule
\end{tabular}\end{center}

\gkpart{解答题}
\begin{itemize}[leftmargin=2.6em]
  \item[\textbf{16.}] $\theta=\arccos\dfrac{\sqrt5}{3}$。
  \item[\textbf{17.}] 一个有意义的逆向问题是“点 $P(2,1)$ 到直线 $3x+4y=t$ 的距离为 $2$，求实数 $t$”；其解为 $t=0$ 或 $t=20$。
  \item[\textbf{18.}] (1) $\dfrac{x^2}{4}+\dfrac{y^2}{2}=1$；(2) $S_{\triangle F_1BN}=\dfrac83$
  \item[\textbf{19.}] (1) 四边形 $EFGH$ 为正方形；(2) $BE=0.3$ 米、$EC=CF=0.1$ 米、$DF=0.3$ 米时，材料费用最省
  \item[\textbf{20.}] (1) $AB=\sqrt6+\sqrt2$；(2) $a^2+b^2<4R^2$；(3) 在 $0<b\le a$、$R>0$ 下，不存在三角形当且仅当 $a>2R$ 或 $a=b=2R$；存在一个三角形当且仅当 $0<b=a<2R$ 或 $0<b<a=2R$；存在两个三角形当且仅当 $0<b<a<2R$。存在时 $c=\dfrac{a\sqrt{4R^2-b^2}\pm b\sqrt{4R^2-a^2}}{2R}$；$b<a<2R$ 时两种符号均取，$a=b<2R$ 时只取加号，$b<a=2R$ 时两式重合为 $\sqrt{4R^2-b^2}$
  \item[\textbf{21.}] (1) $B_1+B_2+\cdots+B_n=nq+\dfrac{n(n-1)}2$；(2) $c_1+c_3-2c_2=q^2>0$，所以 $c_1+c_3>2c_2$；(3) 仅选择第①问。$q=-\dfrac12$，$m$ 只有一个可能值，即 $m=3$，此时前三项的公比为 $\dfrac32$
\end{itemize}
